\chapter{Software Architecture and Engineering Design}
\label{ch:architecture}

\begin{chapterabstract}
This chapter translates the methodology of Chapter~\ref{ch:methods} into a
concrete software system. It presents the layered architecture, a
component-by-component breakdown of the monitor and scheduler, the per-frame
control flow, the policy abstraction that makes baselines and the proposed
scheduler interchangeable, the hyperparameter configuration, and the
hardware-execution strategy used for evaluation.
\end{chapterabstract}

\section{System Architecture Overview}
\label{sec:arch-overview}

The implementation is organised as a modular Python package in which each
stage of Figure~\ref{fig:arch} is a self-contained class with a narrow
interface. A driver iterates over frames, invokes the always-on monitor, and
delegates the compute decision to a pluggable \emph{policy} object. This
separation lets every baseline and the proposed scheduler share identical
perception and evaluation code, so that measured differences are attributable
solely to the allocation decision. The package layout is summarised in
Table~\ref{tab:modules}.

\begin{table}[htbp]
\centering
\caption[Source module layout]{Source module layout and responsibilities.}
\label{tab:modules}
\small
\begin{tabular}{@{}ll@{}}
\toprule
\textbf{Module} & \textbf{Responsibility} \\
\midrule
\texttt{detection/}      & Frame filtering and YOLOv8-n object detection \\
\texttt{tracking/}       & ByteTrack identity association wrapper \\
\texttt{embeddings/}     & CLIP ViT-B/32 embedding extraction \\
\texttt{semantic\_memory/} & Temporal memory and drift estimation \\
\texttt{scheduler/}      & Semantic cache and tiered VLM scheduler \\
\texttt{vlm/}            & Moondream wrapper and VLM factory \\
\texttt{policies/}       & The five interchangeable allocation policies \\
\texttt{evaluation/}     & Metric computation and SCE evaluator \\
\texttt{observatory/}    & Semantic-evolution analytics, visualisation, dashboard \\
\texttt{robotics/}       & Downstream perception-layer adapter \\
\bottomrule
\end{tabular}
\end{table}

\section{Component-Level Implementation Breakdown}
\label{sec:arch-components}

\subsection{Frame Filtering and Detection}
\label{sec:arch-detection}

The frame filter is a CPU-only first guard that discards frames offering no
opportunity for change --- for example near-duplicate frames under a cheap
pixel-difference test --- before any accelerator work is scheduled. Surviving
frames are passed to the nano-scale YOLOv8 detector~\cite{Jocher2023YOLOv8},
which returns class labels and bounding boxes. Only the class--count summary
$O_t$ and the boxes (for the tracker) are retained; pixels are not held beyond
the current iteration, keeping the memory footprint bounded.

\subsection{Tracking and Embedding}
\label{sec:arch-track-embed}

The tracking wrapper feeds detector boxes to ByteTrack~\cite{Zhang2022ByteTrack}
and exposes the set of active identities $\mathcal{I}_t$. In parallel, the
embedding extractor computes the $\ell_2$-normalised CLIP
embedding~\cite{Radford2021CLIP} $E_t$. Both run within the Tier-1 budget and
together with detection constitute the entire cost of \emph{deciding} whether
to reason.

\subsection{Temporal Memory and Drift Estimation}
\label{sec:arch-drift}

The \texttt{semantic\_memory} module implements
Equations~(\ref{eq:memory})--(\ref{eq:dtrack}). It owns the decayed memory
$M_t$, the reference object set $O_{\text{ref}}$, and the reference identity
set $\mathcal{I}_{\text{ref}}$, exposing a single \texttt{drift(frame\_state)}
call that returns the scalar $D_t$ and its three components for logging.
A monitoring helper accumulates drift over time so that slow, silent semantic
drift --- change too gradual to cross a threshold on any single frame --- can be
detected and surfaced. Listing~\ref{lst:drift} shows the core computation.

\begin{lstlisting}[language=Python, caption={Composite drift computation (excerpt).}, label={lst:drift}]
def drift(self, E_t, O_t, ids_t):
    # embedding drift against decayed memory (unit-norm vectors)
    d_embed = 0.5 * (1.0 - float(E_t @ self.M))
    # object-set novelty (Jaccard over class-count pairs)
    inter = self._multiset_intersect(O_t, self.O_ref)
    union = self._multiset_union(O_t, self.O_ref)
    d_obj = 1.0 - (inter / union) if union else 0.0
    # tracking-identity churn (symmetric difference)
    sym = len(ids_t ^ self.ids_ref)
    uni = len(ids_t | self.ids_ref)
    d_trk = (sym / uni) if uni else 0.0
    D_t = self.a * d_embed + self.b * d_obj + self.g * d_trk
    self._update_memory(E_t)            # M <- normalize(lambda*M + (1-lambda)*E_t)
    return D_t, (d_embed, d_obj, d_trk)
\end{lstlisting}

\subsection{Scheduler, Cache, and VLM Wrapper}
\label{sec:arch-scheduler}

The \texttt{scheduler} module realises the tier mapping of
Equation~(\ref{eq:policy}). On a Tier-1 frame it returns the cached output; on
Tier-2/3 it calls a VLM wrapper, which abstracts the underlying model behind a
uniform \texttt{caption(frame, mode)} interface so that the backend can be
swapped through a factory (\texttt{vlm\_factory.py}) without touching the
scheduler. The cache stores the latest output and the state that produced it,
and re-anchors the references on Tier-3 invocations.

\subsection{Vision-Language Model Backends}
\label{sec:arch-vlm-backends}

Because the scheduler treats the VLM as a black box, several backends are
supported and can be selected at run time; this also enables a
\emph{hierarchical} configuration in which the cheaper Tier-2 query and the
heavier Tier-3 query are served by different-capacity models. Two families are
implemented --- Moondream2~\cite{Moondream2024} and the
Qwen2.5-VL~\cite{Qwen25VL} series --- and any compact captioning/VQA model
exposing the same interface can be added. Table~\ref{tab:vlm-backends} lists
the supported and compatible backends together with their fit on a single
Kaggle NVIDIA~T4 (16\,GB).

\begin{table}[htbp]
\centering
\caption[Supported VLM backends]{Vision-language backends and their fit on a
single Kaggle T4 (16\,GB). Moondream2 is the lightweight default; Qwen2.5-VL is
the stronger Tier-3 reasoner. The lower block lists additional black-box-%
compatible models.}
\label{tab:vlm-backends}
\small
\begin{tabular}{@{}llll@{}}
\toprule
\textbf{Backend} & \textbf{Params} & \textbf{T4 (16\,GB) fit} & \textbf{Typical tier role} \\
\midrule
Moondream2                & $\sim$1.9\,B & $\sim$4\,GB (trivial)        & Tier 2 (lightweight, default) \\
Qwen2.5-VL-3B             & $\sim$3\,B   & $\sim$7\,GB (fp16)          & Tier 2/3 (mid) \\
Qwen2.5-VL-7B             & $\sim$7\,B   & $\sim$14--16\,GB (fp16)     & Tier 3 (full) \\
Qwen2.5-VL-7B (4-bit)     & $\sim$7\,B   & $\sim$9\,GB (NF4)           & Tier 3 (full, recommended) \\
\midrule
\multicolumn{4}{@{}l}{\itshape Additional compatible backends:} \\
SmolVLM-2.2B              & $\sim$2.2\,B & $\sim$5\,GB                 & Tier 2 \\
InternVL2-2B / 4B         & 2--4\,B      & $\sim$6--9\,GB              & Tier 2/3 \\
LLaVA-1.5-7B (4-bit)      & $\sim$7\,B   & $\sim$8\,GB                 & Tier 3 \\
\bottomrule
\end{tabular}
\end{table}

In the reported experiments Moondream2 serves the lightweight tier for its low
latency and memory footprint, while Qwen2.5-VL provides stronger full
reasoning; on a two-GPU Kaggle session the heavier model can be pinned to a
second T4 so that loading it never disturbs the always-on Tier-1 monitor.
Crucially, the scheduling behaviour --- call rate, cache hit rate, and SCE ---
is governed by the drift policy and is therefore independent of which backend
is chosen; swapping the model changes the \emph{quality} of each invoked
caption, not \emph{when} the VLM is invoked.

\section{Control Flow and the Policy Abstraction}
\label{sec:arch-controlflow}

Each baseline in Chapter~\ref{ch:experiments} is expressed as a policy object
implementing one method, \texttt{decide(drift, frame\_idx) -> tier}. This
makes the five strategies --- every-frame, uniform sampling, motion gating,
embedding-threshold, and the proposed fused-drift scheduler --- drop-in
replacements that share all other code. The driver loop is:

\begin{enumerate}
  \item read and (optionally) filter the next frame;
  \item run detection, tracking, and embedding to obtain $(O_t,\mathcal{I}_t,E_t)$;
  \item compute $D_t$ via the drift estimator;
  \item ask the active policy for a tier;
  \item execute the tier (reuse cache, light VLM, or full VLM);
  \item record per-frame metrics and update memory.
\end{enumerate}

\section{Hyperparameter Configuration}
\label{sec:arch-hyperparams}

All tunable quantities are declared in YAML so that experiments are
reproducible and the only difference between runs is configuration.
Table~\ref{tab:hyperparams} lists the defaults used unless a study states
otherwise.

\begin{table}[htbp]
\centering
\caption[Default hyperparameters]{Default hyperparameter configuration.}
\label{tab:hyperparams}
\small
\begin{tabular}{@{}lll@{}}
\toprule
\textbf{Symbol} & \textbf{Meaning} & \textbf{Default} \\
\midrule
$\alpha,\beta,\gamma$ & Drift component weights & $0.5,\,0.3,\,0.2$ \\
$\tau_{\text{low}}$   & Tier-1/2 threshold      & $0.15$ \\
$\tau_{\text{high}}$  & Tier-2/3 threshold      & $0.40$ \\
$\lambda$             & Memory decay factor     & $0.90$ \\
$d$                   & CLIP embedding dim.\    & $512$ (ViT-B/32) \\
---                   & Detector                & YOLOv8-n (COCO) \\
---                   & VLM (Tier 2 / Tier 3)   & Moondream2 / Qwen2.5-VL \\
\bottomrule
\end{tabular}
\end{table}

\section{The Semantic Observatory Subsystem}
\label{sec:arch-observatory}

The Semantic Observatory of Section~\ref{sec:meth-observatory} is realised as
the \texttt{observatory} package and the \texttt{run\_observatory.py} driver.
It operates \emph{post hoc} on the \texttt{frame\_log.csv} emitted by any
pipeline run, so it imposes no cost on the real-time path and can analyse the
proposed scheduler or any baseline identically. Its components are:

\begin{itemize}
  \item \texttt{analytics.py} --- the \texttt{SemanticAnalyticsEngine}, which
  computes the ten per-frame signals of
  Table~\ref{tab:observatory-signals} and detects stability windows,
  transition zones, object-persistence records, and semantic events.
  \item \texttt{visualizer.py} --- renders time-aligned figures (drift,
  velocity, acceleration, volatility, tier, and VLM-invocation timelines on a
  shared axis).
  \item \texttt{output\_writer.py} --- serialises the analysis to CSV/JSON
  artefacts for the experiment record.
  \item \texttt{dashboard.py} --- assembles an interactive HTML dashboard for
  qualitative inspection of a run.
  \item \texttt{live\_collector.py} --- an optional hook to accumulate the same
  signals online during a run.
\end{itemize}

Because the Observatory only reads logs and never writes to the scheduler, the
experimental protocol of Equation~(\ref{eq:protocol}) is kept honest: analysis
and decision-making are strictly separated, and every reported improvement is
traceable to a single, recorded change.

\section{Hardware Optimization and Execution Strategy}
\label{sec:arch-hardware}

The system is built on PyTorch~\cite{Paszke2019PyTorch} and
OpenCV~\cite{Bradski2000OpenCV}. Experiments run on a Kaggle NVIDIA~T4 GPU for
GPU-bound stages (CLIP and the VLM) and a commodity laptop CPU for the
filtering and tracking stages, mirroring the heterogeneous compute of a real
edge robot. The engineering levers that keep the Tier-1 monitor within
real-time budget are: (i) running detection and embedding at reduced input
resolution; (ii) keeping the VLM resident in accelerator memory to avoid
reload latency on Tier-2/3 transitions; and (iii) restricting all per-frame
state to compact summaries ($E_t$, $O_t$, $\mathcal{I}_t$) rather than raw
frames, which bounds peak memory irrespective of stream length.
